% Generated draft: requires mathematical and source review.
\documentclass[11pt]{article}
\usepackage[margin=25mm]{geometry}
\usepackage{fontspec}
\setmainfont{Latin Modern Roman}
\usepackage{amsmath,amssymb,mathtools,unicode-math}
\setmathfont{Latin Modern Math}
\usepackage{enumitem,needspace,xurl,hyperref}
\setlength{\parindent}{0pt}
\setlength{\parskip}{5pt}
\setlength{\emergencystretch}{3em}
\newcommand{\N}{\mathbb{N}}
\newcommand{\Z}{\mathbb{Z}}
\newcommand{\Q}{\mathbb{Q}}
\newcommand{\R}{\mathbb{R}}
\newcommand{\C}{\mathbb{C}}
\newcommand{\F}{\mathbb{F}}
\newcommand{\A}{\mathbb{A}}
\newcommand{\PP}{\mathbb{P}}
\newcommand{\E}{\mathbb{E}}
\newcommand{\eps}{\varepsilon}
\newcommand{\dd}{\,\mathrm{d}}
\begin{document}
\section*{UnsolvedMath ID: 30004691}
\textbf{Source problem code:} catalogue entry
\subsection{Definitions and mathematical statement}
Let \(n\geq 1\) be an integer and let \(\Omega\subset\C^n\) be a domain, meaning a nonempty connected open set. For a precise weak formulation, write a bounded \((0,1)\)-form as
\[
f=\sum_{j=1}^n f_j\,\dd\bar z_j,\qquad f_j\in L^\infty(\Omega),
\]
and set
\[
\|f\|_{L^\infty(\Omega)}=\operatorname*{ess\,sup}_{z\in\Omega}\left(\sum_{j=1}^n|f_j(z)|^2\right)^{1/2}.
\]
Interpret \(\bar\partial f=0\) distributionally, that is,
\[
\frac{\partial f_j}{\partial\bar z_k}=\frac{\partial f_k}{\partial\bar z_j}\quad\text{in distributions on }\Omega\quad(1\leq j,k\leq n).
\]
For a fixed domain \(\Omega\), the uniform bounded solvability property asked about is
\[
\begin{gathered}
\exists\,C_\Omega\in(0,\infty)\ \ \forall\,f\in L^\infty_{(0,1)}(\Omega)\text{ with }\bar\partial f=0,\\
\exists\,u\in L^\infty(\Omega):\qquad
\frac{\partial u}{\partial\bar z_j}=f_j\ \text{in distributions }(1\leq j\leq n),\qquad
\|u\|_{L^\infty(\Omega)}\leq C_\Omega\|f\|_{L^\infty(\Omega)}.
\end{gathered}
\]
Thus the constant is independent of \(f\); neither uniqueness nor a linear choice of solution is requested. No boundary values are prescribed.

\textbf{Unresolved scope.} The supplied statement asks whether this property holds for a domain \(\Omega\subset\C^n\), without imposing boundedness of \(\Omega\), pseudoconvexity, boundary regularity, finite type, or a restriction on \(n\). Here ``bounded'' explicitly modifies the form \(f\). However, the supplied background describes the intended general problem as concerning arbitrary pseudoconvex finite-type domains in \(\C^n\) for \(n\geq3\). These formulations have different hypotheses. The record does not specify the intended boundary regularity or meaning of finite type, or whether the domains in that narrower formulation must be bounded. Consequently, the display specifies the analytic property precisely, but the class of domains over which it is being asked remains unresolved. The measurable-coefficient and distributional formulation above is an explicit convention; the source does not specify the regularity of \(f\) or the sense of either equation.
\subsection{Short English statement}
The question asks whether every bounded, \(\bar\partial\)-closed \((0,1)\)-form on the intended domain has a bounded solution of \(\bar\partial u=f\), with one estimate constant depending only on the domain. The statement leaves the domain unrestricted, while the background points to pseudoconvex finite-type domains in dimension at least three. The intended hypotheses and solution regularity require review.
\subsection{Sources}
\begin{itemize}
\item Local UnsolvedMath registry, \texttt{data/unsolvedmath/problems.json}, permanent ID \(30004691\), title ``Uniform Bounded Solutions of the d-Bar Equation,'' source problem code \texttt{catalogue entry}; collection: Oberwolfach Reports Open Problems. The supplied record lists proposed year 2021 and names no proposer.
\item Supplied report citation: \emph{Geometric Methods of Complex Analysis (2022)}, Oberwolfach Reports, DOI 10.4171/owr/2021/25, pp.\ 1291--1345. \url{https://doi.org/10.4171/owr/2021/25}. The record describes these reports as documenting mathematical workshops and their open-problem discussions. The citation's 2022 designation and the record's 2021 proposed year are retained as supplied.
\item Supplied literature source, labelled ``Grundmeier--Simon--Stens\"ones'': \url{https://arxiv.org/abs/1909.04080}.
\item Supplied literature source, labelled ``Yuan, uniform estimates'': \url{https://doi.org/10.1007/s42543-025-00096-0}.
\item Supplied literature assessment dated \textbf{2026-08-22}: the record labels the problem open and states that uniform sup-norm solvability is known for ``\(C^2\) finite-type domains'' and broad classes including real-analytic boundaries, while the general problem for arbitrary pseudoconvex finite-type domains in \(\C^n\), \(n\geq3\), is not fully resolved. This is supplied dated curation derived from the OWR material, not an independently verified literature assessment. The source itself requests independent verification of time-sensitive status and cited claims; its shorthand and omitted hypotheses require review.
\item The supplied statement status is \texttt{corrected\_verified}; its accompanying explanation limits the corrections to formula delimiters, notation, grammar, and typography, with mathematical content and scope preserved. This metadata does not establish mathematical correctness or resolve the scope discrepancy.
\end{itemize}
\end{document}
